\section{Theoretical Framework: The Kronig-Penney Model}
To demonstrate the absence of free electrons in an ideal insulator, we solve the time-independent Schrödinger equation:
$$-\frac{\hbar^2}{2m}\frac{d^2\psi}{dx^2} + V(x)\psi = E\psi$$
where $V(x)$ is a periodic square-well potential representing the atomic lattice. 

\subsection{Energy Gap Formation}
The solution leads to the transcendental equation:
$$P \frac{\sin(\alpha a)}{\alpha a} + \cos(\alpha a) = \cos(ka)$$
where $P$ is the "scattering power" of the potential barrier. In an ideal insulator, the value of $P$ is sufficiently high, and the energy $E$ falls within the "forbidden" region where no real solution for the wave vector $k$ exists. Consequently, electrons are trapped in the valence states with zero probability of occupying the conduction band at $0K$.
